\documentclass[11pt]{article}

\usepackage[margin=1in]{geometry}
\usepackage[T1]{fontenc}
\usepackage{amsmath,amssymb}
\usepackage{array}
\usepackage{longtable}
\usepackage{hyperref}

\newcommand{\code}[1]{\texttt{\detokenize{#1}}}

\title{Peridynamics With Softening: Model and Code Details}
\author{}
\date{}

\begin{document}

\maketitle
\tableofcontents
\newpage

\section{Purpose of This Document}

This document explains the active C++ model in the repository: the implemented
equations, geometry, boundary conditions, softening and failure laws, numerical
relaxation scheme, output files, and the role of each source file.

The current executable is a one-dimensional tensile peridynamic bar. By default
the left grip is fixed, the right grip is pulled in displacement control, and
the virtual substrate coupling is disabled. Several older two-dimensional,
thermal, and explicit film-substrate-grid ideas remain in the source as
comments, but the active calculation uses material points on a line,
\begin{equation}
    y_i = 0,
    \qquad
    i = 0,\ldots,N-1 .
\end{equation}

\section{Source File Map}

\begin{longtable}{|p{0.34\textwidth}|p{0.58\textwidth}|}
\hline
\textbf{File} & \textbf{Purpose} \\
\hline
\code{main.cpp} &
Defines simulation parameters, allocates arrays, builds geometry and correction
factors, calls the mechanical solver, and frees memory. \\
\hline
\code{Geometry.h} &
Builds the one-dimensional nodal coordinates and the peridynamic neighbor
families. \\
\hline
\code{surface_correction_factors.h} &
Computes the end-effect correction factors used in the bond integration
weights. \\
\hline
\code{matrices.h} &
Allocates, initializes, precomputes, and frees vector and pointer-based fields. \\
\hline
\code{only_mechanical.h} &
Contains the main quasi-static mechanical solve: load increments, boundary
conditions, bond forces, softening, damage, virtual substrate coupling, dynamic
relaxation, energy calculation, and output writing. \\
\hline
\code{DefectUtils.h} &
Builds a node mask for an optional defect region. In the current force loop,
the local defect critical stretch is not applied because that code is commented
out. \\
\hline
\code{DiagnosticsUtils.h} &
Provides helpers for diagnostic-node selection and optional diagnostic traces. \\
\hline
\code{time_settings.h} &
Defines mechanical and thermal time-step constants. Only the mechanical values
are active. \\
\hline
\code{variable_initialization.h} &
Initializes scalar scratch variables used by the included-style C++ code. \\
\hline
\code{write_ovito_files.h} &
Legacy OVITO writer. It is not included by the active \code{main.cpp}. \\
\hline
\code{plot_strain.py} &
Post-processing script for plotting or animating \code{result/conf*.dat}. \\
\hline
\code{Makefile} &
Builds the executable \code{peridynamics_with_softening}. \\
\hline
\end{longtable}

\section{Build and Run}

The executable is built with
\begin{verbatim}
make
\end{verbatim}

The default run is
\begin{verbatim}
./peridynamics_with_softening
\end{verbatim}

The executable also accepts optional command-line arguments:
\begin{verbatim}
./peridynamics_with_softening [increments_to_run] [load_increments] [damage_model]
\end{verbatim}

The first argument controls how many outer increments are executed. The second
argument controls the denominator used in the strain increment,
\begin{equation}
    \Delta\varepsilon =
    \frac{\varepsilon_{\mathrm{target}}}{N_{\mathrm{load}}}.
\end{equation}

The third argument selects the damage law:
\begin{center}
\begin{tabular}{|c|l|}
\hline
\textbf{damage\_model} & \textbf{Law} \\
\hline
\(-1\) & Elastic, no damage \\
\(0\) & Brittle failure \\
\(1\) & Linear softening, default \\
\(2\) & Bilinear softening \\
\(3\) & Cornelissen softening \\
\(4\) & Exponential softening \\
\hline
\end{tabular}
\end{center}

For example,
\begin{verbatim}
./peridynamics_with_softening 40 80 3
\end{verbatim}
runs 40 increments with the load-step size of an 80-increment ramp using the
Cornelissen softening law.

\section{Geometry and Discretization}

The active domain is a one-dimensional bar of length
\begin{equation}
    L = 1 .
\end{equation}

The current number of material points is
\begin{equation}
    N = 300 .
\end{equation}

The grid spacing is
\begin{equation}
    \Delta x = \frac{L}{N-1}.
\end{equation}

The material-point coordinates are
\begin{equation}
    x_i = i\Delta x,
    \qquad
    y_i = 0,
    \qquad
    i=0,\ldots,N-1 .
\end{equation}

The active cross-sectional area is
\begin{equation}
    A = 1,
\end{equation}
so the nodal volume is
\begin{equation}
    V = A\Delta x = \Delta x .
\end{equation}

The C++ variables are:
\begin{center}
\begin{tabular}{|l|l|l|}
\hline
\textbf{Quantity} & \textbf{Variable} & \textbf{Current value} \\
\hline
Length & \code{length} & \code{1} \\
Number of points & \code{ndivx}, \code{totnode}, \code{totint} & \code{300} \\
Area & \code{area_bar} & \code{1} \\
Thickness & \code{thickness} & \code{1}, not active in the force law \\
Maximum family storage & \code{maxfam} & \code{20} \\
\hline
\end{tabular}
\end{center}

\section{Horizon and Peridynamic Families}

The peridynamic horizon is
\begin{equation}
    \delta = 3\Delta x .
\end{equation}

The horizon ratio is therefore
\begin{equation}
    m = \frac{\delta}{\Delta x} = 3 .
\end{equation}

For each material point \(i\), its family is
\begin{equation}
    \mathcal{H}_i =
    \left\{
        j \ne i :
        |x_j-x_i| \le \delta + 10^{-12}
    \right\}.
\end{equation}

The small \(10^{-12}\) tolerance prevents floating-point roundoff from dropping
bonds located exactly at the horizon.

The family data are stored using older one-based indexing conventions:
\begin{itemize}
    \item \code{numfam[i]} stores the number of neighbors of node \(i\);
    \item \code{pointfam[i]} stores the one-based offset into \code{nodefam};
    \item \code{nodefam} stores one-based neighbor node indices.
\end{itemize}

When the code retrieves a neighbor, it converts the stored index back to a
zero-based C++ index.

\section{Material Parameters}

The current active material values are:
\begin{center}
\begin{tabular}{|l|l|l|}
\hline
\textbf{Quantity} & \textbf{Variable} & \textbf{Value} \\
\hline
Poisson ratio & \code{pratio} & \(1/3\) \\
Film density & \code{rho_film} & \(1\) \\
Substrate density & \code{rho_substrate} & \(1\), mostly inactive \\
Film Young's modulus & \code{emod_film} & \(1\) \\
Substrate Young's modulus & \code{emod_substrate} & \(1\), mostly inactive \\
\hline
\end{tabular}
\end{center}

The code computes the film plane-stress modulus
\begin{equation}
    \bar{E}_f =
    \frac{E_f}{1-\nu^2},
\end{equation}
and the film shear modulus
\begin{equation}
    \mu_f =
    \frac{E_f}{2(1+\nu)}.
\end{equation}
The active one-dimensional bond force law uses the peridynamic bond stiffness
below, not the shear modulus directly.

\section{Bond Micromodulus}

The one-dimensional bond stiffness before discrete correction is
\begin{equation}
    c_0 =
    \frac{2E}{A\delta^2}.
\end{equation}

The code applies the correction
\begin{equation}
    C_d = \frac{m}{m+1}.
\end{equation}

Thus the active stiffness is
\begin{equation}
    c =
    \frac{2E}{A\delta^2}
    \frac{m}{m+1}.
\end{equation}

In the code this is stored as \code{bc_mechanical_film_bar}.

\section{Surface and Volume Corrections}

Boundary-near nodes have incomplete horizons. The active geometric correction is
based on the visible left and right horizon lengths:
\begin{equation}
    \delta_i^- = \min(\delta,x_i),
\end{equation}
\begin{equation}
    \delta_i^+ = \min(\delta,L-x_i).
\end{equation}

The nodal geometric correction is
\begin{equation}
    G_i =
    \frac{2\delta}{\delta_i^-+\delta_i^+}.
\end{equation}

The first and last three displacement-prescribed grip nodes are explicitly set
to
\begin{equation}
    G_i = 1 .
\end{equation}

The bond correction is the average
\begin{equation}
    G_{ij} =
    \frac{G_i+G_j}{2}.
\end{equation}

The separate volume-correction factor is currently
\begin{equation}
    \mathrm{fac}_{ij}=1 .
\end{equation}

The effective integration weight for a bar-bar bond is
\begin{equation}
    w_{ij} = V\,G_{ij}\,\mathrm{fac}_{ij}.
\end{equation}

\section{Kinematics}

The displacement of material point \(i\) is
\begin{equation}
    u_i .
\end{equation}

The current position is
\begin{equation}
    y_i = x_i+u_i .
\end{equation}

For bond \((i,j)\), the reference separation is
\begin{equation}
    \xi_{ij} = x_j-x_i,
\end{equation}
with length
\begin{equation}
    r_{ij}^0 = |\xi_{ij}|.
\end{equation}

The current separation is
\begin{equation}
    \eta_{ij}
    =
    (x_j+u_j)-(x_i+u_i),
\end{equation}
with length
\begin{equation}
    r_{ij}=|\eta_{ij}|.
\end{equation}

The bond stretch is
\begin{equation}
    s_{ij}
    =
    \frac{r_{ij}-r_{ij}^0}{r_{ij}^0}.
\end{equation}

The one-dimensional direction cosine is
\begin{equation}
    n_{ij} =
    \frac{\eta_{ij}}{r_{ij}}.
\end{equation}

The active implementation computes \(r_{ij}\) using the \(x\)-component only,
which is consistent with the active one-dimensional geometry.

For history-dependent damage, the code stores
\begin{equation}
    s_{ij}^{\max}(t)
    =
    \max_{\tau\le t} s_{ij}(\tau).
\end{equation}
This field is \code{bond_s_history[i][j]}.

\section{Load Control and Boundary Conditions}

The imposed macroscopic strain starts from
\begin{equation}
    \varepsilon_0=0 .
\end{equation}

The default target strain is
\begin{equation}
    \varepsilon_{\mathrm{target}} = 0.004.
\end{equation}

The default number of load increments is
\begin{equation}
    N_{\mathrm{load}} = 80.
\end{equation}

The strain increment is
\begin{equation}
    \Delta\varepsilon =
    \frac{\varepsilon_{\mathrm{target}}}{N_{\mathrm{load}}}
    =
    5\times 10^{-5}
    \quad
    \text{for the defaults.}
\end{equation}

At outer load step \(k\),
\begin{equation}
    \varepsilon_k =
    \varepsilon_{k-1}+\Delta\varepsilon.
\end{equation}

\subsection{Load Predictor}

Before dynamic relaxation begins, the code keeps the previous converged
displacement field and adds an affine predictor:
\begin{equation}
    u_i
    \leftarrow
    u_i + 0.75\,\Delta\varepsilon\,x_i .
\end{equation}

\subsection{Displacement-Prescribed Grip Nodes}

The first and last three nodes are prescribed:
\begin{equation}
    i=0,1,2,
    \qquad
    i=N-3,N-2,N-1.
\end{equation}

During every inner relaxation step, these nodes are reset to
\begin{equation}
    u_i = 0,
    \qquad i=0,1,2,
\end{equation}
and
\begin{equation}
    u_i = \varepsilon_k L,
    \qquad i=N-3,N-2,N-1.
\end{equation}
The prescribed-node velocities are reset to
\begin{equation}
    v_i = 0 .
\end{equation}

Therefore the left boundary satisfies
\begin{equation}
    u(0)=0,
\end{equation}
and the right boundary satisfies
\begin{equation}
    u(L)=\varepsilon_k L.
\end{equation}

The active force loop skips the prescribed nodes, setting their damage, strain
energy, substrate energy, and stress to zero.

\subsection{Body Forces}

The body-force array is initialized to zero and no active body force is applied:
\begin{equation}
    b_i=0 .
\end{equation}

The old traction-style loading code is present only as comments.

\section{Virtual Substrate Coupling}

The virtual substrate code path is now optional and disabled by default:
\begin{equation}
    \code{enable_substrate_coupling}=\code{false},
    \qquad
    \code{c_int_ratio}=0.
\end{equation}
Therefore, in the simpler tensile-bar scenario,
\begin{equation}
    p_i^{\mathrm{sub}}=0,
    \qquad
    W_i^{\mathrm{sub}}=0.
\end{equation}

If the flag is re-enabled in the source, the substrate is represented by a
virtual affine layer, not a separate deformable mesh. It is located a vertical distance
\begin{equation}
    h=\Delta x
\end{equation}
below the bar.

The virtual substrate displacement is prescribed as
\begin{equation}
    u_j^{\mathrm{sub}} =
    \varepsilon_k x_j .
\end{equation}

For bar node \(i\) and virtual substrate point \(j\), the initial interface
distance is
\begin{equation}
    r_{ij}^{0,\mathrm{sub}}
    =
    \sqrt{(x_j-x_i)^2+h^2}.
\end{equation}

The interaction is included only when
\begin{equation}
    r_{ij}^{0,\mathrm{sub}}
    \le
    \delta + 10^{-12}.
\end{equation}

The current horizontal separation is
\begin{equation}
    d_{ij}^{\mathrm{sub}}
    =
    (x_j+u_j^{\mathrm{sub}})
    -
    (x_i+u_i).
\end{equation}

The current interface distance is
\begin{equation}
    r_{ij}^{\mathrm{sub}}
    =
    \sqrt{
        \left(d_{ij}^{\mathrm{sub}}\right)^2+h^2
    }.
\end{equation}

The interface stretch is
\begin{equation}
    s_{ij}^{\mathrm{sub}}
    =
    \frac{
        r_{ij}^{\mathrm{sub}}-r_{ij}^{0,\mathrm{sub}}
    }{
        r_{ij}^{0,\mathrm{sub}}
    }.
\end{equation}

The horizontal direction cosine is
\begin{equation}
    n_{ij}^{\mathrm{sub}}
    =
    \frac{d_{ij}^{\mathrm{sub}}}{r_{ij}^{\mathrm{sub}}}.
\end{equation}

The interface stiffness is
\begin{equation}
    c_{\mathrm{int}} =
    \alpha_{\mathrm{int}}c,
    \qquad
    \alpha_{\mathrm{int}}=1 .
\end{equation}

When enabled, the virtual substrate force contribution is
\begin{equation}
    p_i^{\mathrm{sub}}
    =
    \sum_j
    c_{\mathrm{int}}
    s_{ij}^{\mathrm{sub}}
    n_{ij}^{\mathrm{sub}}
    V .
\end{equation}

The optional substrate sum uses \(V\) directly. It does not use the bar-bar surface
correction \(G_{ij}\).

\section{Bar-Bar Bond Force}

For an intact elastic bond, the force contribution from \(j\) to \(i\) is
\begin{equation}
    f_{ij}
    =
    c\,s_{ij}\,n_{ij}\,w_{ij}.
\end{equation}

The bar-bar internal force is
\begin{equation}
    p_i^{\mathrm{bar}}
    =
    \sum_{j\in\mathcal{H}_i}
    f_{ij}.
\end{equation}

The total force used by dynamic relaxation is
\begin{equation}
    p_i =
    p_i^{\mathrm{bar}}
    +
    p_i^{\mathrm{sub}}.
\end{equation}

\section{Damage and Softening Models}

The active default is
\begin{equation}
    \code{degradation_model}=1,
\end{equation}
which means linear softening.

The implemented options are:
\begin{center}
\begin{tabular}{|c|l|}
\hline
\textbf{Value} & \textbf{Meaning} \\
\hline
\(-1\) & Elastic, no damage \\
\(0\) & Brittle failure \\
\(1\) & Linear softening \\
\(2\) & Bilinear softening \\
\(3\) & Cornelissen softening \\
\(4\) & Exponential softening \\
other & Fallback to brittle-style behavior \\
\hline
\end{tabular}
\end{center}

\subsection{Code-to-Equation Mapping in the Bond Loop}

For each directed bond from node \(i\) to neighbor \(j\), the active code first
computes the kinematic quantities
\begin{align}
    r_{ij}^0 &= \code{idist}, \\
    w_{ij} &= \code{w_const}
           = V\,G_{ij}\,\mathrm{fac}_{ij}, \\
    n_{ij} &= \code{dirx}
           = \frac{\eta_{ij}}{|\eta_{ij}|}, \\
    s_{ij} &= \code{stretch}
           = \frac{|\eta_{ij}|-r_{ij}^0}{r_{ij}^0}.
\end{align}

The historical maximum stretch stored in \code{bond_s_history[i][j]} is updated
before the softening branch is evaluated:
\begin{equation}
    \hat{s}_{ij}^{\,n+1}
    =
    \max\left(
        \hat{s}_{ij}^{\,n},
        s_{ij}^{\,n+1}
    \right).
\end{equation}
Below, \(\hat{s}_{ij}\) denotes this updated history value.

The code also defines a loading-envelope flag:
\begin{equation}
    \mathrm{onEnvelope}_{ij}
    =
    \left(
        |s_{ij}-\hat{s}_{ij}| < 10^{-14}
    \right).
\end{equation}
When this is true, the current stretch is equal to the historical maximum to
roundoff tolerance, so the bond is on the active softening envelope. When it is
false, the bond is unloading or reloading below a previously reached maximum.

The force contribution added to node \(i\) is the scalar
\begin{equation}
    \Delta p_{ij}
    =
    \code{dforce1_mechanical}.
\end{equation}
The bar-bar force accumulator uses
\begin{equation}
    p_i^{\mathrm{bar}}
    \leftarrow
    p_i^{\mathrm{bar}}+\Delta p_{ij}.
\end{equation}

The stress moment accumulator uses the same force contribution:
\begin{equation}
    M_i
    \leftarrow
    M_i+\Delta p_{ij}\,\xi_{ij},
\end{equation}
and the reported stress is
\begin{equation}
    \sigma_i = \frac{1}{2}M_i.
\end{equation}

For nodal damage, the code accumulates two local sums:
\begin{align}
    A_i &\leftarrow A_i + q_{ij}V\,\mathrm{fac}_{ij}, \\
    B_i &\leftarrow B_i + V\,\mathrm{fac}_{ij},
\end{align}
where \(q_{ij}=1\) for intact elastic bonds, \(q_{ij}=\mu_{ij}\) for softened
bonds, and \(q_{ij}=0\) for fully failed bonds. The final nodal damage is
\begin{equation}
    D_i = 1-\frac{A_i}{B_i}.
\end{equation}

Because the family table stores directed bonds, the code also checks the
reciprocal entry \((j,i)\). If either \(\mathrm{fail}_{ij}\) or the reciprocal
entry has failed, the bond is treated as failed and carries no force.

If \(\mathrm{fail}_{ij}=0\) when the bond loop reaches the bond, the bond is
already fully failed and the active code uses
\begin{align}
    \Delta p_{ij} &= 0, \\
    W_{ij} &= 0, \\
    \Delta G_{ij}^{\mathrm{fracture}} &= 0, \\
    q_{ij} &= 0.
\end{align}
That is, already failed bonds continue to count in the denominator \(B_i\), but
they do not add force, strain energy, or damage numerator.

\subsection{Common Elastic Contribution}

Whenever a branch decides that a bond is elastic, the equations entering the
code are
\begin{align}
    \Delta p_{ij}
    &=
    c\,s_{ij}\,w_{ij}\,n_{ij}, \\
    W_{ij}
    &=
    \frac{1}{4}
    c\,s_{ij}^{2}\,r_{ij}^0\,w_{ij}, \\
    \Delta G_{ij}^{\mathrm{fracture}}
    &=
    0, \\
    q_{ij}
    &=
    1 .
\end{align}

\subsection{Common Softening Contribution}

For the linear, bilinear, Cornelissen, and exponential softening branches, the
irreversibility variable is the history stretch
\begin{equation}
    \kappa_{ij}(t)
    =
    \max_{\tau\le t}s_{ij}^{+}(\tau),
\end{equation}
stored in the code as \code{bond_s_history[i][j]}. The prescribed envelope is
written as
\begin{equation}
    f_{\mathrm{env},ij}(\kappa)
    =
    c\,p(\kappa).
\end{equation}
For the implemented softening laws, the code has
\(p(\kappa)=s_0\mu_{ij}\) after the elastic limit. The degraded
unloading/reloading stiffness is
\begin{equation}
    g(\kappa)=\frac{p(\kappa)}{\kappa},
    \qquad
    k_{\mathrm{unload},ij}
    =
    c\,g(\kappa)
    =
    \frac{f_{\mathrm{env},ij}}{\kappa}.
\end{equation}

If the bond is on the envelope,
\begin{equation}
    \Delta p_{ij}
    =
    f_{\mathrm{env},ij}\,w_{ij}\,n_{ij}.
\end{equation}

If the bond is unloading or reloading,
\begin{equation}
    \Delta p_{ij}
    =
    k_{\mathrm{unload},ij}
    s_{ij}
    w_{ij}
    n_{ij}.
\end{equation}

In both cases, the strain-energy contribution stored by the code is the
directed-bond version of the free energy:
\begin{equation}
    W_{ij}
    =
    \frac{1}{4}
    c\,g(\kappa_{ij})
    s_{ij}^{2}
    r_{ij}^0
    w_{ij}.
\end{equation}
The factor \(1/4\) appears because the implementation accumulates directed
bond contributions at nodes; it is the pair-split counterpart of the usual
\(\frac12 c r g(\kappa)s^2\) bond free energy.

The damage numerator uses
\begin{equation}
    q_{ij}=\mu_{ij}.
\end{equation}

\subsection{Common Complete-Failure Contribution}

For the brittle, linear, bilinear, Cornelissen, and fallback branches, once the
code decides that a bond has completely failed, it sets
\begin{equation}
    \mathrm{fail}_{ij}=0,
\end{equation}
and uses
\begin{align}
    \Delta p_{ij} &= 0, \\
    W_{ij} &= 0, \\
    q_{ij} &= 0.
\end{align}

For the brittle, bilinear, Cornelissen, exponential, and fallback branches, the
current code adds the branch-local complete-failure value stored in
\code{d_fracture_energy}. For \(\code{degradation_model}=1\), the code instead
uses the incremental history-energy update described in the linear-softening
subsection. For linear softening at complete failure, the limiting accumulated
value is
\begin{equation}
    G_{ij}^{\mathrm{fracture}}
    =
    \frac14 c\,s_0s_c\,r_{ij}^0w_{ij}.
\end{equation}

\subsection{Critical Stretch}

The initial critical stretch is
\begin{equation}
    s_0 =
    \sqrt{
        \frac{3G_c}{E\delta}
    }.
\end{equation}

The current fracture parameter is
\begin{equation}
    G_c = 10^{-8},
\end{equation}
and the softening energy parameter is
\begin{equation}
    R_d = G_c.
\end{equation}

The complete-failure stretch \(s_c\) depends on the selected model.

For linear softening,
\begin{equation}
    s_c =
    1.5
    \left(
        s_0+
        \frac{2R_d}{cs_0}
    \right).
\end{equation}

For bilinear softening,
\begin{equation}
    s_c =
    s_0+
    \frac{2R_d}{cs_0}.
\end{equation}

For the Cornelissen law,
\begin{equation}
    s_c =
    s_0+
    \frac{5.1361R_d}{cs_0}.
\end{equation}

For the exponential law,
\begin{equation}
    s_c = 10s_0
\end{equation}
is used as a numerical reference cutoff in \code{main.cpp}; the active
exponential branch now uses this value as the complete-failure cutoff.

The bilinear transition stretch is
\begin{equation}
    s_k =
    s_0 + 0.15(s_c-s_0),
\end{equation}
with residual ratio
\begin{equation}
    \beta = 0.25 .
\end{equation}

\subsection{Elastic Model}

For \(\code{degradation_model}<0\), the bond remains elastic:
\begin{equation}
    f_{ij}
    =
    c\,s_{ij}\,n_{ij}\,w_{ij}.
\end{equation}

\subsection{Brittle Model}

For \(\code{degradation_model}=0\), the bond is elastic while
\begin{equation}
    s_{ij} \le s_0 .
\end{equation}

If
\begin{equation}
    s_{ij} > s_0,
\end{equation}
then
\begin{equation}
    f_{ij}=0,
    \qquad
    \mathrm{fail}_{ij}=0 .
\end{equation}

The fracture-energy accumulator receives
\begin{equation}
    G_{ij}^{\mathrm{fracture}}
    =
    \frac{1}{4}
    c\,s_0^2\,r_{ij}^0\,w_{ij}.
\end{equation}

\subsection{Linear Softening Model}

If
\begin{equation}
    s_{ij}^{\max} \le s_0
\end{equation}
the bond is elastic.

If
\begin{equation}
    s_0 < s_{ij}^{\max} < s_c,
\end{equation}
the degradation factor is
\begin{equation}
    \mu_{ij}
    =
    \frac{s_c-s_{ij}^{\max}}{s_c-s_0}.
\end{equation}

The envelope force magnitude is
\begin{equation}
    f_{\mathrm{env}}
    =
    cs_0\mu_{ij}.
\end{equation}

On the loading envelope,
\begin{equation}
    f_{ij}
    =
    f_{\mathrm{env}}n_{ij}w_{ij}.
\end{equation}

For unloading or reloading below the historical maximum, the secant stiffness is
\begin{equation}
    k_{\mathrm{unload}}
    =
    \frac{f_{\mathrm{env}}}{s_{ij}^{\max}},
\end{equation}
and
\begin{equation}
    f_{ij}
    =
    k_{\mathrm{unload}}s_{ij}n_{ij}w_{ij}.
\end{equation}

If
\begin{equation}
    s_{ij}^{\max}\ge s_c,
\end{equation}
the bond is fully failed:
\begin{equation}
    f_{ij}=0,
    \qquad
    \mathrm{fail}_{ij}=0 .
\end{equation}

For \(\code{degradation_model}=1\), fracture energy is stored incrementally with
a per-bond history value \(\code{bond_fracture_energy_history[i][j]}\). During
partial softening the code computes
\begin{equation}
    W_{\mathrm{input}}
    =
    \frac12 c s_0^2
    +
    \frac12 c s_0(1+\mu_{ij})(s_{ij}^{\max}-s_0),
\end{equation}
\begin{equation}
    W_{\mathrm{recoverable}}
    =
    \frac12 f_hs_{ij}^{\max},
    \qquad
    f_h=cs_0\mu_{ij},
\end{equation}
and
\begin{equation}
    G_{\mathrm{now}}
    =
    \frac12
    \left(
        W_{\mathrm{input}}-W_{\mathrm{recoverable}}
    \right)
    r_{ij}^0w_{ij}.
\end{equation}
Only the positive increment
\begin{equation}
    \Delta G
    =
    G_{\mathrm{now}}-G_{\mathrm{history}}
\end{equation}
is added to \(\code{fracture_energy[i]}\), and then \(G_{\mathrm{history}}\) is
updated. At complete failure, the total target value is
\begin{equation}
    G_{\mathrm{total}}
    =
    \frac14 c s_0s_c r_{ij}^0w_{ij}.
\end{equation}

\subsection{Bilinear Softening Model}

For \(s_0<s_{ij}^{\max}<s_k\),
\begin{equation}
    \mu_{ij}
    =
    1
    -
    (1-\beta)
    \frac{s_{ij}^{\max}-s_0}{s_k-s_0}.
\end{equation}

For \(s_k\le s_{ij}^{\max}<s_c\),
\begin{equation}
    \mu_{ij}
    =
    \beta
    \frac{s_c-s_{ij}^{\max}}{s_c-s_k}.
\end{equation}

Then the same envelope and unloading/reloading formulas are used.

\subsection{Cornelissen Softening Model}

Define
\begin{equation}
    \eta =
    \frac{s_{ij}^{\max}-s_0}{s_c-s_0}.
\end{equation}

The degradation factor is
\begin{equation}
    \mu_{ij}
    =
    \left(1+(3\eta)^3\right)e^{-6.93\eta}
    -
    28\eta e^{-6.93}.
\end{equation}

The code clamps this value:
\begin{equation}
    \mu_{ij}
    =
    \max(0,\mu_{ij}).
\end{equation}

\subsection{Exponential Softening Model}

The exponential law uses
\begin{equation}
    \alpha =
    \frac{cs_0}{R_d},
\end{equation}
and
\begin{equation}
    \mu_{ij}
    =
    \exp\left[
        -\alpha
        \left(s_{ij}^{\max}-s_0\right)
    \right].
\end{equation}

If
\begin{equation}
    s_{ij}^{\max}\ge s_c,
\end{equation}
the bond is fully failed:
\begin{equation}
    f_{ij}=0,
    \qquad
    \mathrm{fail}_{ij}=0 .
\end{equation}

\section{Nodal Damage Variable}

The nodal damage-like variable is
\begin{equation}
    D_i
    =
    1
    -
    \frac{
        \sum_{j\in\mathcal{H}_i}
        \mu_{ij}V\,\mathrm{fac}_{ij}
    }{
        \sum_{j\in\mathcal{H}_i}
        V\,\mathrm{fac}_{ij}
    }.
\end{equation}

For intact elastic bonds, \(\mu_{ij}=1\). For softened bonds, the current
history-based \(\mu_{ij}\) is used. Fully failed bonds contribute to the
denominator but not the numerator.

Thus \(D_i=0\) means all counted bonds around node \(i\) are intact; larger
values indicate local softening or bond failure. Grip nodes are assigned
\(D_i=0\) because their force loop is skipped.

\section{Stress Measure}

The stress-like output is a peridynamic force-moment measure:
\begin{equation}
    \sigma_i
    =
    C_\sigma
    \sum_{j\in\mathcal{H}_i}
    f_{ij}\xi_{ij}.
\end{equation}

The active coefficient is
\begin{equation}
    C_\sigma=\frac{1}{2}.
\end{equation}

The optional virtual substrate force is included in equilibrium only when the
source flag
\begin{equation}
    \code{enable_substrate_coupling}
\end{equation}
is true. It is not included in this bar-bar stress moment.

\section{Output Strain}

The saved strain is a finite-difference post-processing field, not the
peridynamic bond stretch.

At the left endpoint,
\begin{equation}
    \varepsilon_0 =
    \frac{u_1-u_0}{\Delta x}.
\end{equation}

For interior nodes,
\begin{equation}
    \varepsilon_i
    =
    \frac{1}{2}
    \left[
        \frac{u_i-u_{i-1}}{\Delta x}
        +
        \frac{u_{i+1}-u_i}{\Delta x}
    \right].
\end{equation}

At the right endpoint,
\begin{equation}
    \varepsilon_{N-1}
    =
    \frac{u_{N-1}-u_{N-2}}{\Delta x}.
\end{equation}

\section{Energy Quantities}

For an elastic intact bond, the strain-energy contribution at node \(i\) is
\begin{equation}
    W_{ij}
    =
    \frac{1}{4}
    c\,s_{ij}^2\,r_{ij}^0\,w_{ij}.
\end{equation}

For softened unloading or reloading, the code uses
\begin{equation}
    W_{ij}
    =
    \frac{1}{4}
    k_{\mathrm{unload}}
    s_{ij}^2
    r_{ij}^0
    w_{ij}.
\end{equation}

For \(\code{degradation_model}=1\), fracture energy is accumulated
incrementally from the bond history. The code stores a previous per-bond value
\(\code{bond_fracture_energy_history[i][j]}\), computes the current history
value \(G_{\mathrm{now}}\), and adds only
\begin{equation}
    \Delta G = G_{\mathrm{now}}-G_{\mathrm{history}}
\end{equation}
when this increment is positive. This prevents the same history energy from
being added repeatedly during ADR iterations. In the current code, the other
damage branches keep their branch-local complete-failure increments.

If substrate coupling is enabled, the virtual substrate energy is
\begin{equation}
    W_i^{\mathrm{sub}}
    =
    \sum_j
    \frac{1}{4}
    c_{\mathrm{int}}
    \left(s_{ij}^{\mathrm{sub}}\right)^2
    r_{ij}^{0,\mathrm{sub}}
    V.
\end{equation}
For the current default tensile-bar scenario, this term is identically zero.

The kinetic energy is
\begin{equation}
    K_i =
    \frac{1}{2}m_i v_i^2.
\end{equation}

The current lumped mass-like value is
\begin{equation}
    m_i=1.
\end{equation}

Damping dissipation is accumulated as
\begin{equation}
    \mathcal{D}_i^{n+1}
    =
    \mathcal{D}_i^n
    +
    m_i c_d v_i^2\Delta t .
\end{equation}

The total stored nodal energy density is
\begin{equation}
    E_i^{\mathrm{tot}}
    =
    W_i
    +
    W_i^{\mathrm{sub}}
    +
    K_i
    +
    \mathcal{D}_i
    +
    G_i^{\mathrm{fracture}}.
\end{equation}

\section{Dynamic Relaxation Solver}

Each outer load increment is equilibrated by an explicit damped dynamic
relaxation loop. The active equation is
\begin{equation}
    m_i\ddot{u}_i
    +
    c_d\dot{u}_i
    =
    p_i+b_i.
\end{equation}

Since \(b_i=0\),
\begin{equation}
    m_i\ddot{u}_i
    +
    c_d\dot{u}_i
    =
    p_i.
\end{equation}

The acceleration used in the code is
\begin{equation}
    a_i^n =
    \frac{
        p_i^n+b_i-c_dv_i^n
    }{
        m_i
    }.
\end{equation}

The active update is
\begin{equation}
    v_i^{n+1}
    =
    v_i^n+a_i^n\Delta t,
\end{equation}
\begin{equation}
    u_i^{n+1}
    =
    u_i^n+v_i^{n+1}\Delta t.
\end{equation}

The mechanical time step is
\begin{equation}
    \Delta t = \mathrm{dt\_mechanical\_bar}.
\end{equation}

It is computed from
\begin{equation}
    \tau =
    \frac{L}{\sqrt{E/\rho}},
\end{equation}
and
\begin{equation}
    \mathrm{dt\_mechanical\_bar}
    =
    \frac{\mathrm{dt\_mechanical}}{\tau}.
\end{equation}

With the current defaults \(L=E=\rho=1\),
\begin{equation}
    \tau=1,
    \qquad
    \mathrm{dt\_mechanical\_bar}
    =
    5\times 10^{-4}.
\end{equation}

The code also computes \code{dt_pd}, \code{dt_lambda}, and \code{dt_bar}, but
the active mechanical update uses \code{dt_mechanical_bar}.

\section{Damping and Convergence}

The default damping coefficient is
\begin{equation}
    c_d = 5.
\end{equation}

This is the active path because
\begin{equation}
    \code{fixed_damping_coefficient}=\code{true}.
\end{equation}

If adaptive damping is enabled, the code estimates
\begin{equation}
    c_{\mathrm{opt}}
    =
    2\sqrt{
        \frac{\mathrm{num}}{\mathrm{den}}
    },
\end{equation}
caps the result, and smooths it as
\begin{equation}
    c_d
    \leftarrow
    0.7c_d+0.3c_{\mathrm{opt}}.
\end{equation}

The relaxation loop stops when
\begin{equation}
    \max_i |v_i|
    \le
    10^{-6}.
\end{equation}

The maximum iteration count is
\begin{equation}
    100000.
\end{equation}

\section{Algorithm Summary}

The initialization sequence is:
\begin{enumerate}
    \item Read constants from \code{main.cpp}, \code{variable_initialization.h},
    and \code{time_settings.h}.
    \item Allocate arrays with \code{matrices}.
    \item Build the one-dimensional grid and family list with
    \code{build_Geometry}.
    \item Compute surface correction factors with
    \code{surface_correction_factors}.
    \item Precompute bond geometry and weights with
    \code{precompute_bond_invariants}.
    \item Initialize displacement and velocity fields to zero.
    \item Set the lumped mass-like coefficient to \(1\) for every node.
    \item Call \code{mechanical}.
\end{enumerate}

For each outer load increment:
\begin{enumerate}
    \item Increase the applied strain by \(\Delta\varepsilon\).
    \item Add the affine predictor \(0.75\Delta\varepsilon x_i\).
    \item Reset velocities and old force state for dynamic relaxation.
    \item Run the inner relaxation loop until convergence.
    \item Write final fields for the increment.
    \item Optionally detect first crack and stop five increments later.
\end{enumerate}

For each dynamic relaxation iteration:
\begin{enumerate}
    \item Reimpose displacement and zero velocity on the first and last three nodes.
    \item Compute bar-bar bond stretches and forces.
    \item Update historical maximum stretch for each bond.
    \item Apply the selected degradation law.
    \item Add virtual substrate/interface force only if substrate coupling is enabled.
    \item Compute nodal damage, stress, strain energy, and substrate energy.
    \item Update damping if adaptive damping is enabled.
    \item Update velocity and displacement explicitly.
    \item Compute kinetic energy, damping dissipation, total energy, and convergence speed.
    \item Compute finite-difference output strain.
    \item Write optional diagnostics or snapshot buffers.
\end{enumerate}

\section{Output Files}

\subsection{Configuration Files}

For every outer increment, the code writes
\begin{equation}
    \code{result/confK.dat}.
\end{equation}

Each row contains
\begin{equation}
    x_i,\qquad u_i,\qquad \varepsilon_i,\qquad D_i.
\end{equation}

\subsection{OVITO Files}

For every increment, the code writes
\begin{equation}
    \code{ovito_files/for_ovito_K.xyz}.
\end{equation}

The columns are
\begin{equation}
    x,\ y,\ u_x,\ u_y,\ x_{\mathrm{def}},\ y_{\mathrm{def}},\
    D,\ \sigma,\ W,\ \varepsilon .
\end{equation}

Because the active model is one-dimensional, \(y\), \(u_y\), and
\(y_{\mathrm{def}}\) are written as zero.

\subsection{Final ADR Field Matrices}

The final converged fields for each increment are written to:
\begin{itemize}
    \item \code{Final_ADR_stress.txt}
    \item \code{Final_ADR_strain.txt}
    \item \code{Final_ADR_strain_energy.txt}
    \item \code{Final_ADR_kinetic_energy.txt}
    \item \code{Final_ADR_fracture_energy.txt}
    \item \code{Final_ADR_damping_dissipation_energy.txt}
    \item \code{Final_ADR_substrate_energy.txt}
\end{itemize}

Rows correspond to outer increments; columns correspond to nodes.

\subsection{Energy Histories}

System-level energy histories during the ADR iterations are written to:
\begin{itemize}
    \item \code{total_energy_density_matrix.txt}
    \item \code{total_kinetic_energy_density_matrix.txt}
    \item \code{total_damping_dissipation_energy_density_matrix.txt}
    \item \code{total_fracture_energy_density_matrix.txt}
    \item \code{total_strain_energy_density_matrix.txt}
    \item \code{total_substrate_energy_density_matrix.txt}
\end{itemize}

Each row is one outer increment; each column is one ADR iteration.

\subsection{Stress Convergence}

When \code{track_stress_convergence} is true, the code writes
\begin{equation}
    \code{convergence_of_stress.txt}.
\end{equation}

It tracks \(\sigma\) at the node closest to
\begin{equation}
    x=0.3.
\end{equation}

\subsection{Snapshots and Diagnostics}

The snapshot mechanism can write all ADR-step fields for selected increments:
\begin{itemize}
    \item \code{stressK.txt}
    \item \code{strainK.txt}
    \item \code{strain_energyK.txt}
    \item \code{kinetic_energyK.txt}
    \item \code{fracture_energyK.txt}
    \item \code{damping_dissipation_energyK.txt}
    \item \code{substrate_energyK.txt}
    \item \code{total_energy_densityK.txt}
\end{itemize}

By default, \code{write_snaps} is set to a value not reached by ordinary runs.

If \code{track_numerical_noise} is enabled, diagnostic stress, force,
displacement, and maximum-stretch traces are also written for selected nodes.

\section{First-Crack Early Stop}

The early-stop mode is disabled by default:
\begin{equation}
    \code{stop_after_first_crack}=\code{false}.
\end{equation}

When enabled, first crack is detected when any node satisfies
\begin{equation}
    D_i > 0.
\end{equation}

The simulation then stops five increments after that first detection.

\section{Defect Handling}

\code{DefectUtils.h} can mark nodes in a circular defect region:
\begin{equation}
    (x_i-x_c)^2+(y_i-y_c)^2
    \le
    r_{\mathrm{defect}}^2.
\end{equation}

It computes a reduced critical stretch
\begin{equation}
    s_0^{\mathrm{defect}}
    =
    s_0
    \times
    \mathrm{defect\_strength\_factor}.
\end{equation}

The default defect values are:
\begin{center}
\begin{tabular}{|l|l|l|}
\hline
\textbf{Quantity} & \textbf{Variable} & \textbf{Value} \\
\hline
Enabled & \code{introduce_defect} & \code{false} \\
Center & \code{defect_x_center}, \code{defect_y_center} & \((0,0)\) \\
Radius & \code{defect_radius} & \(\Delta x\) \\
Strength factor & \code{defect_strength_factor} & \(0.6\) \\
\hline
\end{tabular}
\end{center}

Important implementation detail: the force loop currently comments out the code
that would replace \(s_0\) with \(s_0^{\mathrm{defect}}\). Therefore enabling
the defect mask alone does not change the active bond failure or softening law.

\section{Post-Processing}

The Python script \code{plot_strain.py} reads \code{result/conf*.dat}. Example
uses:
\begin{verbatim}
python3 plot_strain.py --steps 1 20 80
python3 plot_strain.py --steps 1 20 80 --displacement --damage
python3 plot_strain.py --movie peridynamics.mp4 --damage --strain-zoom
\end{verbatim}

The script expects the columns
\begin{equation}
    x,\qquad u,\qquad \varepsilon,\qquad D.
\end{equation}
If an older file has only three columns, it fills damage with zeros.

\section{Active Versus Dormant Code}

\begin{longtable}{|p{0.35\textwidth}|p{0.57\textwidth}|}
\hline
\textbf{Dormant item} & \textbf{Current status} \\
\hline
Two-dimensional film/substrate grid &
Commented out; the active model uses \(y=0\) for all nodes. \\
\hline
Thermal solve &
Thermal constants exist, but no active thermal loop is called. \\
\hline
Explicit traction boundary conditions &
Present as comments; active loading is fixed-left and pulled-right displacement control. \\
\hline
Separate deformable substrate points &
Disabled by default; optional virtual affine substrate coupling remains in the solver. \\
\hline
\code{write_ovito_files.h} &
Not included by active \code{main.cpp}. \\
\hline
Defect critical stretch &
The mask exists, but the local critical-stretch substitution is commented out. \\
\hline
\code{dt_bar} stability estimate &
Computed for reference, but active ADR uses \code{dt_mechanical_bar}. \\
\hline
Half-step ADR variant &
Older half-step velocity update code is commented out. \\
\hline
\code{heaviside} helper &
Defined in \code{main.cpp}, but unused. \\
\hline
\end{longtable}

\section{Important Implementation Notes}

\begin{enumerate}
    \item The active model is mostly nondimensional and uses variables with the
    suffix \code{_bar}.
    \item The pair \((i,j)\) and the pair \((j,i)\) have separate entries in
    \code{fail} and \code{bond_s_history}.
    \item The stress output includes only bar-bar bond force moments, not the
    optional virtual substrate force.
    \item Grip nodes are skipped by the force loop; damage can occur in the
    free bar bonds once the selected law reaches its damage threshold.
    \item Grip nodes are excluded from the active force loop and report zero
    stress and damage from that path.
    \item \code{maxfam} must be large enough for every family. With \(m=3\) in
    one dimension, \code{maxfam=20} is comfortably sufficient.
    \item Generated files in \code{result/}, \code{ovito_files/}, and the energy
    text files are run products rather than source files.
\end{enumerate}

\end{document}
